\begin{center}\begin{tikzpicture}[font=\small,>=stealth,line join=round]
\draw[gray!40,rounded corners=3pt] (0,0) rectangle (14.2,8.3);\node at (7.1,7.700000000000001) {例3：相邻两项一上一下，函数值相加为0};
\begin{scope}[shift={(1.5918032786885248,3.7966666666666664)},x=0.9836065573770493cm,y=0.30666666666666664cm]\draw[->,gray] (-0.5,0)--(5.6,0) node[right]{$x$};\draw[->,gray] (0,-7)--(0,8) node[above]{$y$};\draw[blue,thick,solid,domain=-0.2:5.3,samples=180] plot(\x,{5*cos(deg(pi*\x))});\fill[orange!85!black] (0,5) circle(2.2pt);\fill[orange!85!black] (1,-5) circle(2.2pt);\fill[orange!85!black] (2,5) circle(2.2pt);\fill[orange!85!black] (3,-5) circle(2.2pt);\fill[orange!85!black] (4,5) circle(2.2pt);\fill[orange!85!black] (5,-5) circle(2.2pt);\foreach \v in {0,1,2,3,4,5}{\ifnum\v=0\node[below left=5pt] at (\v,0){$\v$};\else\node[below=5pt] at (\v,0){$\v$};\fi}\end{scope}\node[anchor=west,align=left,text width=5.0cm] at (8.4,4.0) {$f=5\cos\pi x,\ a=1$\\[9pt]偶数位置：5\\[9pt]奇数位置：$-5$\\[9pt]每两项相消\\[9pt]2026项为1013对};
\node[text width=12.6cm,align=center] at (7.1,.65) {图示前3对；任意满足条件的函数同样可以两项一组相消。};
\end{tikzpicture}\end{center}
